\scanpage{21}
Thus we have the relations
\[
(13)\quad\left\{\begin{gathered}
\tilde\sigma_0^2=\tilde\sigma_1^2=\tilde\sigma_\infty^2=1\\
\tilde\sigma_0\tilde\sigma_1=\tilde\sigma_1\tilde\sigma_\infty=\tilde\sigma_\infty\tilde\sigma_0
=\overset{\deleted{\sim}}{\rho}\ \deleted{(\unsure{=\rho})}\\
(\overset{\deleted{\sim}}{\rho}{}^3=1\text{ is a consequence of the preceding ones})
\end{gathered}\right.
\]
\margin{NB
\[
\rho\tilde\sigma_i\rho^{-1}=\tilde\sigma_{\dot\rho i}\quad\text{i.e}
\]
\[
\begin{aligned}
\rho\tilde\sigma_0\rho^{-1}&=\tilde\sigma_1\\
\rho\tilde\sigma_1\rho^{-1}&=\tilde\sigma_\infty\\
\rho\tilde\sigma_\infty\rho^{-1}&=\tilde\sigma_0
\end{aligned}
\]}
and the inverse image $\tilde{\mathcal E}$ of $\varphi(\overset{\deleted{\unsure{\sim}}}{\unsure{\mathfrak S_3|\mathcal S_3}})\subset\Gamma$ (subgroup of index 2) is a $\frac12$ direct product of $\unsure{\mathfrak S_3|\tilde{\mathfrak S}_3|\mathcal S_3|\tilde{\mathcal S}_3}$ and $\Pi$. We must make explicit the actions of the $\tilde\sigma_a$ on\unsure{.} $\unsure{\pi|\Pi}$ to have the complete structure of the group, We find that $\tilde\sigma_a$ lies over $\varphi(\overset{\deleted{\sim}}{\dot\sigma}_a)$ i.e. the reflection with respect to the plane passing through $P,\bar P,\deleted{\unsure{a_i|\infty|a}}i$, which sends $l_0$ to $l_0^{-1}$, $\deleted{l}l_\infty$ to $l_{\unsure{1|0}}^{-1}$, $l_1$ to $l_\infty^{-1}$ :
\margin{These are those obtained by transport of structure, via the actions of $\unsure{\mathfrak S_3|\tilde{\mathfrak S}_3|\mathcal S_3|\tilde{\mathcal S}_3}$ on $(U,P)$}
\[
(14)\quad\left\{\begin{array}{lll}
\tilde\sigma_0l_0\tilde\sigma_0^{-1}=l_0^{-1}&
\tilde\sigma_0l_1\tilde\sigma_0^{-1}=l_\infty^{-1}&
\tilde\sigma_0l_\infty\tilde\sigma_0^{-1}=l_1^{-1}\\
\tilde\sigma_1l_1\tilde\sigma_1^{-1}=l_1^{-1}&\begin{tikzpicture}[baseline=-0.4ex,x=1em,y=1ex]
\fill (0,0.12)--(0.08,0.32)--(0.15,0.11)--(0.31,0.06)--(0.16,-0.05)--(0.13,-0.22)--(0.04,-0.08)--(-0.09,-0.06)--cycle;
\draw[line width=0.55pt] (0.50,0)--(0.70,0) (0.98,0)--(1.18,0) (1.46,0)--(1.66,0);
\end{tikzpicture}&\\
\begin{tikzpicture}[baseline=-0.4ex,x=1em,y=1ex]
\draw[line width=0.55pt] (0,0)--(0.20,0) (0.48,0)--(0.68,0) (0.96,0)--(1.16,0) (1.44,0)--(1.64,0);
\end{tikzpicture}&&
\end{array}\right.
\]
\margin{\uline{Critique} The notation $\tilde\sigma_0$ rather suggests that $\tilde\sigma_0$ lies over $\sigma_0$ --- it is better to write $\dot\sigma_i'=\dot\sigma_i\tau$ and use the notation $\tilde\sigma_i'$ instead of $\tilde\sigma_i$ $(i\in\{0,1,\infty\})$, or better $\sigma_i'$ (without $\sim$) --- to reserve $\tilde\sigma_i'$ for elements of $S\mathcal E$, cf below\ldots}
The relations (13) (14), plus the fundamental relation $l_\infty l_1l_0=1$ (1), give a presentation of $\tilde{\mathcal E}$, from which one can if one wishes eliminate $\rho$, \unsure{(}by omitting to write the relation which defines it)

We can now lift $\dot\tau$ in three ways, \unsure{namely|namely}
\[
(15)\quad\left\{\begin{aligned}
\tau_0&=\sigma_0\tilde\sigma_0=\tilde\sigma_0\sigma_0\\
\tau_1&=\sigma_1\tilde\sigma_1=\tilde\sigma_1\sigma_1\\
\tau_\infty&=\sigma_\infty\tilde\sigma_\infty\equiv\tilde\sigma_\infty\sigma_\infty
\end{aligned}\right.
\qquad\text{thus}\quad(16)\left\{\begin{aligned}
\tau_0^2&=1\\\tau_1^2&=1\\\tau_\infty^2&=1
\end{aligned}\right.
\]
\mspage{20}
